\documentclass[11pt]{article}
\usepackage[margin=1.03in]{geometry}
\usepackage[expansion=false]{microtype}
\usepackage{amsmath,amssymb,amsthm,mathtools}
\usepackage[colorlinks=true,linkcolor=blue,urlcolor=blue,citecolor=blue]{hyperref}
\newtheorem{theorem}{Theorem}[section]
\newtheorem{lemma}[theorem]{Lemma}
\newtheorem{proposition}[theorem]{Proposition}
\newcommand{\phys}{\mathrm{phys}}
\newcommand{\WF}{\operatorname{WF}}
\newcommand{\supp}{\operatorname{supp}}
\newcommand{\ran}{\operatorname{Ran}}
\newcommand{\vol}{\mathrm{dvol}}
\newcommand{\dd}{\mathrm{d}}
\newcommand{\STF}{\mathrm{TF}}
\DeclareMathOperator{\arctanh}{arctanh}
\title{Local curvature information and modular dynamics\\in coupled Einstein--scalar theory}
\author{}
\date{Research manuscript 09+21, Edition 2, 25 November 2026}
\begin{document}
\maketitle

\begin{abstract}
We construct positive quasifree Hadamard states for the physical
linearized theory of canonical Einstein gravity with scalar fields about smooth
globally hyperbolic four-dimensional solutions, with finitely many
scalars and an arbitrary smooth potential. No spatial compactness,
uniform end geometry, scalar rank, or Killing-field assumption is made.
Local gauge recovery and a smooth signed charge correction give the
exact physical commutator, including after reduction by the compact
presymplectic radical. A single smooth Gaussian preparation makes
every compact physical Ward test lie in all powers of the local modular
logarithm in every region. The bare generator has a convergent
physical Gram reconstruction; relatively compact regions admit a
redundant coordinate distribution. On the original closed four-scalar
round solution, one relational scalar together with curvature generates
the full local algebra; two independently prepared coherent shifts
hide from the curvature restriction, with finite two-sided
relative-entropy bounds after Gaussian preparation. The round example
retains its explicit material-coordinate modular kernel.
\end{abstract}

\section{Introduction}
A local physical state and a local modular generator are different
objects. The first requires a positive covariance with the exact
gauge-reduced causal bracket, while the second also depends on the
chosen state and the local algebra. Neither follows from a formal
frequency splitting of a gauge-fixed wave operator. The compact
critical state construction of \cite{p09}, the flat Weyl entropy
calculation of \cite{p21}, and the general linear gauge frameworks of
\cite{hack,hackschenkel,khavpeierls} motivate the present comparison.

Sections~\ref{sec:general}--\ref{sec:construction} construct a physical
state on a general fixed background. In particular the actual
indefinite wave charge is used when the potential gradient makes
trace reversal fail to commute with evolution. The noncompact step
uses regional gauge recovery and puts no uniform Fredholm assumption
at infinity. Section~\ref{sec:generalmod} prepares one reference
simultaneously for all regions, separates local standardness from
global cyclicity, and establishes a domain-controlled bare modular
reconstruction. The later round worked case retains its distinct
local observable classification, preparations, entropy experiment
and explicit coordinate kernel. Those round claims are not asserted
on arbitrary backgrounds; conversely the general state proof does
not use the scalar markers of that example.

All constructions concern linearized physical observables. We
neither construct a nonlinear material-time Hamiltonian nor assert
interacting constraint quantization, ultraviolet completeness, or
an invariant preferred vacuum.

\section{Canonical gauge complex and the physical theorem}\label{sec:general}
Let $(M,g,\phi)$ be a smooth time-oriented globally hyperbolic
four-dimensional solution of
\[
 S=\tfrac12\int_M\{R-2\Lambda-\textstyle\sum_A|d\phi^A|^2-2V(\phi)\}\vol_g .
\]
A constant invertible field redefinition absorbs any fixed positive
definite target kinetic matrix. The potential $V$ is an arbitrary smooth
function, not assumed critical or bounded below. We consider the real
linear theory on this fixed background; the spaces of configurations and
sources below are those of the linear theory only. The symbols $C_c^\infty$
and $H^s_{\rm loc}$ are used for sections of the relevant smooth real or
complex bundles. We use complexification for covariance and symbol
calculations. Adjoints denoted $*$ use auxiliary positive pairings unless
specified; $\dagger$ denotes the action adjoint.

\begin{theorem}[Positive physical reference]\label{thm:generalstate}
For every such background and every smooth Cauchy surface, the compact
Ward/equation phase space has a regular positive quasifree state with the
exact causal physical commutator and signed Hadamard two-point lifts.
The state may be chosen to descend through the algebraic presymplectic
radical. No condition on the topology, spatial ends, scalar rank,
Killing fields, or uniform geometry of the Cauchy surface is required.
The signed lifts are wave-gauge bisolutions on the redundant field
bundle; their commutator is prescribed on physical sources, not on
arbitrary redundant sources.
\end{theorem}

\subsection{Wave identities and noncritical charge}
Write $X\mapsto KX=(\mathcal L_X g,X\phi)$ for gauge. In the original
metric variable $h$, set $Ih=h-\frac12g\operatorname{tr}_g h$ and
\[
 b_2((h,\psi),(h',\psi'))=\tfrac14\bar h:Ih'+\bar\psi\cdot\psi',
 \qquad
 F(h,\psi)_a=\nabla^b h_{ba}-\tfrac12\nabla_a\operatorname{tr}h
       -2\psi^A\nabla_a\phi^A .
\]
The ghost pairing $b_1$ is the Lorentzian one. Integration by parts gives
$K^\dagger=-F/2$. Let $A_{ab}=\sum_A a_{Aa}a_{Ab}$, where $a_A=d\phi^A$, and
$W=\Lambda+V$. In these variables the normalized Euler Hessian is
$P(h,\psi)=(-2\delta(\operatorname{Ric}-A-Wg),
\delta(\Box\phi-V'))$. In the first component the Euler trace
reversal has been moved into $b_2$; the variation of that trace
reversal on the background equation is zero. On shell the exact identities
are $PK=0$, $FP=0$, and
\[
 D_2=P+KF,\qquad D_1=FK,\qquad
 D_2K=KD_1,\quad FD_2=D_1F .
\]
Both $D_i$ are normally hyperbolic and formally self-adjoint for their
respective $b_i$. To see the cancellation relevant to the scope, put
$H_A=\nabla^2\phi^A$. The non-diagonal zeroth-order
blocks of $D_2$ are $D_{h\psi}^A=-4H_A+2V_Ag$ and
$(D_{\psi h}h)^A=-H_A:h$; the scalar diagonal potential is
$-V_{AB}-2a_A\cdot a_B$. The metric diagonal is
$\Box h_{ab}+2R_{acbd}h^{cd}-2A_{(a}{}^c h_{b)c}$ in the
convention $g^{cd}R_{acbd}=\operatorname{Ric}_{ab}$, and
$D_1=\Box+\operatorname{Ric}-2A$ on ghosts. In the metric row
the $a_{A(a}\nabla_{b)}\psi^A$ variation of $A$ cancels that
from the symmetrized gradient of $-2\psi^Aa_A$ in $KF$.
In the scalar row the differentiated metric appears through
$-a_A\cdot(\operatorname{div} h-\tfrac12d\operatorname{tr}h)$,
which cancels against the scalar row of $KF$. The on-shell identities are
$\operatorname{Ric}=\sum_A a_A\otimes a_A+(\Lambda+V)g$ and
$\operatorname{tr} H_A=V_A$. In particular
$I(-4H_A+2V_Ag)=-4H_A$, but
\[
 [D_2,\operatorname{diag}(I,1)](h,\psi)
       =(2V_A\psi^Ag,\tfrac12V_A\operatorname{tr}h).
\]
It is the actual indefinite charge, not a separately conserved
unreversed tensor charge, that is used here. Compare the underlying
gauge complex in \cite{hack,hackschenkel,schmid}. The existence
results for Maxwell fields in \cite{schmid} and for vacuum
linearized gravity on asymptotically flat backgrounds in \cite{gmw}
are not used as state-existence theorems for this coupled system.

Let $U_i$ be Cauchy evolution, $\rho_i$ the value and first normal
derivative trace on $\Sigma$, and $E_i$ the causal Green operator,
with the Peierls sign chosen consistently throughout. Let $Q_i$ denote
the Hermitian charge $i$ times the sesquilinear wave symplectic form
with integrand $b_i(u_0,v_1)-b_i(u_1,v_0)$. Fix the sign of the
positive sheet by $u_1=-i|\xi|u_0$ in a flat normal chart. Conservation
follows from formal $b_i$-self-adjointness and the parallel principal
connection. Put $k=\rho_2KU_1$, $c=\rho_1FU_2$, and $Q=Q_2$.
On smooth wave cut data,
\begin{equation}
 ck=0,\qquad Q(ka,v)=-\tfrac12Q_1(a,cv).                 \label{eq:gaugecharge}
\end{equation}
The second identity follows from the Green pairing and $E_2K=KE_1$;
one may first test with $v=\rho_2E_2 f$, which realizes every compact
smooth cut datum by a spacelike compact Cauchy solution and a time
switch. It extends by compact-test duality.

The real phase space below is
$\mathcal V=\ker F|_{C_c^\infty}/P C_c^\infty$, with its real Peierls
form $\sigma$. It is canonically the compact data in $\ker c$ modulo
compact $k$-gauge data. For surjectivity realize a compact constrained
wave datum by $E_2 f$; then $E_1Ff=0$, so $Ff=D_1a$ for compact $a$,
and $f-Ka$ is Ward. For injectivity, if $E_2 f=KX$ with compact ghost
cut data, take $X=E_1b$ and write $f=Kb+D_2u$ with compact $b,u$.
Then $D_1(b+Fu)=0$ and compact homogeneous uniqueness gives $f=Pu$.
These arguments use the exact normally hyperbolic Green sequences.
The image of a region always means its compact source image in
$\mathcal V$, not an independently reduced regional quotient.

\subsection{ADM comparison, including an off-constraint identity}
Write $M_\Sigma$ for the differential map from wave cut variables
to ADM metric and scalar pairs, and $S_\Sigma$ for the following
differential section. Set the perturbed lapse and shift to zero on
the cut, fix the spatial metric and normal derivatives from the
prescribed canonical pairs, and solve the four remaining normal
derivatives of the lapse/shift from $F|_\Sigma=0$. This is algebraic
with invertible normal de Donder coefficient. It gives
$M_\Sigma S_\Sigma=1$ even off constraint. If the ADM constraints
hold, the normal Euler equations and $D_2u=0$ give
$\partial_nF|_\Sigma=0$: all tangential derivatives of $F$ vanish
there, and the normal projections of the metric row of $KF$ have
invertible coefficients on the four remaining derivatives. Thus
$F=0$ propagates by $D_1$. Conversely propagated $F=0$ implies the
ADM constraints. The same statements hold distributionally.

If $w$ is a compact constrained wave datum then
$w-S_\Sigma M_\Sigma w=k(0,a_1)$ for compact $a_1$: a zero-value
ghost with suitable normal derivative changes precisely the
lapse/shift values without affecting ADM pairs. The remaining normal
derivatives match by $F|_\Sigma=0$. The ADM gauge operator $B$ acts
on the lapse and shift descriptors; $M_\Sigma k$ is $B$ of the
ghost value, and $S_\Sigma B$ is compact wave gauge on compact
descriptors (choose its ghost normal derivative to keep lapse/shift
zero). Denote the real ADM charge by $Q_a$, with the factor of $i$
understood when using its Hermitian complexification.

The charge identity for two $S_\Sigma$ lifts needs only
$F|_\Sigma=0$, not the constraints. Indeed the gauge-fixed quadratic
action adds $-b_1(Fu,Fu)/4$ modulo a divergence. Its variation adds
normal boundary potential $-b_1(Fu,\sigma_F(n)\delta u)/2$, up to
the choice of outward normal. Its antisymmetrized polarization uses
the values of $F$ only. The ungauged action current equals the
canonical ADM current modulo a tangential divergence by the lapse--shift
form of the same action; the hypersurface boundary functional contributes
a field-space exact variation and cancels in the antisymmetrization.
Consequently $Q(S_\Sigma a,S_\Sigma b)=Q_a(a,b)$ when one argument
is compact smooth, including a distributional first argument. In
addition, (\ref{eq:gaugecharge}) pairs the gauge gap $k(0,a_1)$
of a constrained datum to zero against $S_\Sigma a$, even when $a$
is off constraint, because the first value of $cS_\Sigma a$ is
zero. In particular, for every compact Ward $f$ and any cut
distribution $a$,
\begin{equation}
 Q(S_\Sigma a,\rho_2E_2 f)=Q_a(a,M_\Sigma\rho_2E_2 f).       \label{eq:offcharge}
\end{equation}
The Cauchy shadow of $\supp f$ is compact, so the distributional Green
pairing in (2) only uses a fixed compact Cauchy buffer; no integration
at spatial infinity is implicit.

\section{Construction of the positive state}\label{sec:construction}
\subsection{Proper microlocal positive squares}
Use balanced Cauchy symbols of orders $H^{1/2}\oplus H^{-1/2}$,
i.e. locally conjugate by elliptic symbols of orders $1/2,-1/2$.
No global bounded balancing inverse is used. The real-principal wave
system has roots $\pm|\xi|$ with their graph eigenbundles. Successively
eliminate each off-diagonal degree by division by the nonzero root
difference, and Borel sum with proper kernels. This gives projectors
$C_{j,s}$, $s=\pm1$, modulo properly smoothing operators, whose sum
may be taken exactly one. Their wave lifts are the one-sheet Cauchy
Fourier integral propagators \cite{gow}. The constructions are local
in space on compact cut charts and require no uniform spatial constants.
For the exact wave intertwiners $k,c$, the opposite-sheet blocks are
smoothing to every degree: in their evolution equations the first
possible symbol is multiplied by the root difference one order above
its time derivative. The same argument, using conservation
$Q'=i(QA-A^*Q)$ for $i\partial_t=A$ on the evolving data, gives
$(C_{2,-s})^*Q C_{2,s}\equiv0$. Here and below $\equiv$ is equality
modulo proper all-order smoothing; hence $C_{2,s}^*Q\equiv QC_{2,s}$.

The symbol of $k$ on each branch at null $\zeta$ is $(\zeta\odot X,0)$;
it is injective because $\zeta(n)\ne0$, and the two signs are
independent. That of $c$ is surjective: for target $f$ use
$Ih=\ell\otimes f+f\otimes\ell-f(\zeta^\sharp)\ell\otimes\ell$
with $\ell(\zeta^\sharp)=1$. Choose a proper left symbol parametrix
$L_0$ of $k$ using $k^*k$, and proper right parametrix $R_0$ of $c$.
Set
\[
 L=\sum_s C_{1,s}L_0C_{2,s},\quad T=1-kL,\quad
 R=T\sum_s C_{2,s}R_0 C_{1,s},\quad H=1-Rc .
\]
Thus $cT=c$ and $ck=0$ exactly, $Tk\equiv0$, $T^2\equiv T$,
$cR\equiv1$, $TR\equiv R$, and $HT\equiv TH$. On $\ker c$,
$H=1$ exactly and $T$ changes inputs only by gauge. Define
$e_s=C_{2,s}TH$. It is an oblique projector modulo smoothing.
On its principal range, choose temporal gauge $h_{n\mu}=0$ and
use $\zeta^a Ih_{ab}=0$ to get two spatial TT polarizations; the
scalar positions are unaffected at principal order. The signed
charge there is proportional to
$|h_{\rm TT}|^2/4+|\psi|^2>0$ on each nonzero physical representative.
The principal gauge shift costs zero charge by
(\ref{eq:gaugecharge}), so the result also holds in the chosen
representative. The full matrix
$e_s^*(sQ)e_s+(1-e_s)^*(1-e_s)$ then has a strictly positive
principal symbol. Starting with its positive principal root, solve
each subleading coefficient by the invertible Sylvester map
$X\mapsto b_0X+Xb_0$ and properly sum to obtain $b_s$ whose
square $b_s^*b_s$ has this full symbol. Thus the actual factors
$F_s=b_s e_s$ satisfy
\begin{equation}
 F_s^*F_s\equiv sH^*T^*Q C_{2,s}TH,qquad F_s k\equiv0 .       \label{eq:formal}
\end{equation}
Indeed $e_s^*Qe_s\equiv H^*T^*Q C_{2,s}TH$ by the charge adjoint
and idempotence, and $(1-e_s)e_s\equiv0$. No positive low-frequency
spectrum of an indefinite wave form is invoked. Both $H$ factors in
(3) are essential before restricting to constraints.

\subsection{Exact regional gauge annihilation}\label{sec:localgauge}
We describe the local recovery used for this purpose. The symmetric
gradient in the metric row of $KX$ determines all derivatives of the prolonged
Killing pair $(X,\nabla_{[a}X_{b]})$ up to a connection term and
known differential expressions in the metric and scalar sources.
In a coordinate trivialization write this prolongation as
$(d+\mathsf A)s=F_1(u)$, with $u=KX$ and $s$ the pair, and
include the algebraic residuals of the prolongation and scalar row.
On a sufficiently small star-shaped ball use the regularized
Poincar\'e homotopy $J$ of order $-1$ \cite{costabel}. For $1$-forms
it maps into degree zero, and on degree zero $Jd=1-P_0$ where
$P_0$ is a smooth average taking values among constants. The inverse
of $1+J\mathsf A$ exists by the small-radius $L^2$ norm bound for
$J\mathsf A$, and builds a particular recovery
$s_0=(1+J\mathsf A)^{-1}JF_1(u)$. The remaining homogeneous
candidates are the finite-dimensional smooth family
$(1+J\mathsf A)^{-1}c$ for constants $c$. Fit their residuals in
the simultaneous prolonged and algebraic equations by a fixed
finite set of compact smooth duals to their residual image; directions
of zero residual are left untouched. For a pure gauge input this
recovers the exact pair modulo a joint Killing solution. For general
inputs the same construction is linear and the residual
$G(u)=u-K\Xi u$ satisfies $GK=0$.

Here are the support and regularity facts used from this construction.
On nested buffered balls, the transpose of each homotopy maps
compact smooth tests into compact smooth tests, with support
contained in the convex hull with its fixed interior center bump.
For the Neumann inverse the same holds: the iterates for the transpose
remain in that fixed compact hull, and the inverse equation gives
successively one Sobolev order by the smoothing order of $J^*$.
Localized elliptic regularity for $1+J\mathsf A$ extends recovery to
distributions by transpose and shows pseudolocality. The fitted
finite-dimensional terms are smooth and use compact smooth duals.
In particular $G^*$ of a compact test gives an actual compact Ward
source, after the identification of auxiliary and action pairings.
These arguments use the finite-type Killing prolongation and make no
constant-rank assumption on the residual image as the ball varies.

On a buffered cut patch, apply $\Xi$ to the local $D_2$ solution
and take the recovered ghost value and first normal derivative.
Call this cut map $L_i$. For exact pure gauge input in that buffer,
the recovered ghost differs from the actual ghost by a local joint
Killing solution, so $1-kL_i$ annihilates pure gauge there. Here $k$
uses the wave-defined ghost Cauchy map, not arbitrary second
derivatives of $\Xi u$. The map and its transpose preserve compact
smooth data on the indicated buffers: for the transpose the inserted
slice test is a compact timelike-conormal distribution (delta and its
normal derivative). Pseudolocal recovery keeps the wavefront set inside
this noncharacteristic conormal, and the \emph{causal} propagator converts it to smooth
cut data by the Green identity, with buffered support by finite
propagation. Compact buffers may be chosen locally finite.

Choose $\sum_i\chi_i^2=1$ and $\eta_i=1$ near $\supp\chi_i$ with
support inside the working buffer of $L_i$. Define
\[
 D_{i,s}=\chi_i F_s\eta_i(1-kL_i),\qquad
 \lambda_s^0=\sum_i D_{i,s}^*D_{i,s}.
\]
Each factor kills $k$ exactly. The difference from $\chi_i F_s$
is smoothing: separated supports handle the missing $1-\eta_i$;
commuting $\eta_i$ through differential $k$ uses $F_sk\equiv0$
and a separated-support commutator. An outer cutoff within the
working buffer allows extension by zero. The bounded smooth transpose
mapping just established controls the remaining composition in both
kernel arguments. The sum is locally finite, hence
$\lambda_s^0\equiv F_s^*F_s$. It is positive on the full wave cut
and has the correct signed wavefront after the two Cauchy propagations
from (3); their common diagonal cut covector excludes a one-sided
zero covector, including at caustics.

\subsection{Closed constraint projection and the smooth signed defect}
On the full wave cut the Hermitian form
$R_w=H^*T^*QTH-(\lambda_+^0-\lambda_-^0)$ is jointly smooth
by (3). Pull only this smooth defect to ADM:
$R_a=S_\Sigma^*R_wS_\Sigma$. Choose a positive ungraded ADM
$L^2$ space $\mathcal H$ with smooth weight $w\geq1$ growing as
necessary so that $R_a$ represents a bounded Hilbert--Schmidt operator
and every compactly localized output derivative is Hilbert--Schmidt.
For example, bound the kernel and derivatives on successive compact
shell pairs and multiply each argument by $w^{-1/2}$ so the resulting
squared shell norms are summable; for each fixed derivative order
enforce these bounds from that shell onward. Localization makes
derivatives of the fixed weight harmless. The real charge is bounded
and has pointwise invertible smooth Riesz matrix $J_w$ (no bounded
global inverse is asserted).

The distributional constraint subspace is
$\mathcal C=(\overline{J_w B C_c^\infty})^{\perp_{\mathcal H}}$.
Its orthogonal projector is locally regular. Here are the local
estimates and their origin. On bounded nested coordinate domains,
the shift metric gauge row controls the symmetric gradient of the
shift, up to zero-order lapse; the lapse momentum row controls its
Hessian minus its trace, up to first-order shift and lower lapse.
The pointwise Hessian-minus-trace map is invertible on symmetric
tensors in spatial dimension three. By ordinary local Korn and
Hessian estimates one obtains, for $D=J_wB$,
\[
 \|\alpha\|_{H^2}+\|\beta\|_{H^1}
       \le C(\|D(\alpha,\beta)\|_{L^2}+\|\alpha\|_{L^2}+\|\beta\|_{L^2})
\]
on smooth bounded cut subdomains. The first-order shift term in
the second row is controlled using the first row and interpolation
absorbs first derivatives of the lapse. On the orthogonal complement
of the local kernel the two $L^2$ terms can be discarded by Rellich
and a contradiction sequence. Thus an element of the global gauge
range closure has a local potential with bounded graded norm, by
projecting an approximating sequence off the local kernel on a
slightly larger domain. If $v$ is the Hilbert projection of $f$ into
that range closure, this potential satisfies $D^*D a=D^*f$ locally.
The mixed elliptic normal parametrix and bounded local potential give
\begin{equation}
 \|P_{\mathcal C}f\|_{H^m(U)}
 \le C_{m,U,U'}(\|f\|_{H^m(U')}+\|f\|_{\mathcal H}),
 \qquad U\Subset U'.                                      \label{eq:localell}
\end{equation}
All weights here are bounded above and below locally; no global Korn
or global Fredholm estimate has been used.

For the radical-reduced conclusion let $\mathcal C_0$ be the
$\mathcal H$ closure of compact smooth constrained vectors, and
$\mathcal E=\mathcal C\cap\mathcal C_0^\perp$. We prove
(\ref{eq:localell}) also for $P_{\mathcal C_0}$. Fix
$e\in\mathcal E$ and put $a=J_w^{-1}e$ locally distributionally.
Its charge pairs to zero with all compact constrained ADM vectors.
By (\ref{eq:offcharge}), the distributional wave solution
$U_2 S_\Sigma a$ annihilates every compact Ward test. In particular
it annihilates all the regional $G^*$ tests of
Section~\ref{sec:localgauge}, hence is locally $K\Xi (U_2 S_\Sigma a)$.
The pure gauge parameter has well-defined cut traces: the wave
solution has only characteristic wavefront and the recovery is
pseudolocal on the smaller buffer. The ADM readout then gives
$a=B\xi$ locally. Since $e\in\mathcal C$,
$(J_wB)^*(J_wB)\xi=0$ distributionally. The mixed normal
operator is elliptic as above, so $e$ is smooth. For a fixed compact
localization the map $\mathcal E\to H^m_{\rm loc}$ has a closed
graph by distributional convergence, hence its local seminorm is
bounded by the $\mathcal H$ norm. Therefore
$P_{\mathcal C_0}=P_{\mathcal C}-P_{\mathcal E}$ obeys (4).
This proof neither identifies $\mathcal C_0$ with $\mathcal C$
nor assumes a finite-dimensional end space.

Compress $R_a$ by $P_{\mathcal C_0}$ to get $R_0^a$.
By (4) its compactly localized output derivatives are
Hilbert--Schmidt: apply (4) to the Hilbert--Schmidt map with both
its global $\mathcal H$ norm and its localized derivative norm.
Right multiplication by a bounded orthogonal projection preserves
this. Selfadjointness gives the analogous input estimates; Fourier
weights localized in the two variables then give jointly arbitrary
Sobolev order, hence a smooth kernel. The same proof for the positive
and negative spectral parts uses
$|R_0^a|=R_0^a\operatorname{sgn}(R_0^a)$, so local output
bounds are retained before using selfadjointness.

The defect $R_a$ pairs compact gauge to zero against all
distributionally constrained vectors in $\mathcal H$:
$H=1$ on both arguments, $T$ shifts by gauge, and exact compact
Ward integration gives this for the charge term; the positive
factors annihilate gauge identically. Thus $R_0^a$ and its signs
kill compact gauge. If $r$ is a compact algebraic radical,
(2) and the regional $G^*$ argument give a locally pure gauge
smooth wave lift $S_\Sigma r$. Since it is constrained, its
local ghost solves $D_1$ by $FK=D_1$. Hence all $D_{i,s}$ kill
$S_\Sigma r$. By definition of the radical $R_a(r,f)=0$ for all
compact constrained $f$, and by boundedness for the closure
$\mathcal C_0$. It follows that $R_0^a$ and its signs kill $r$.
Only the implication radical $\Rightarrow$ locally gauge is used;
no global Killing-cocycle classification is needed.

Finally set
\[
 \lambda^s=\lambda_s^0+M_\Sigma^*(R_0^a)_s M_\Sigma.
\]
These are positive. The difference on constrained compact arguments
is exactly $Q$, because $P_{\mathcal C_0}$ fixes their ADM readout,
$S_\Sigma M_\Sigma-1$ is compact gauge there, and $H=1$ there.
They descend through the radical by the preceding argument.
The correction is smooth even on the ambient wave cut since
$M_\Sigma$ is differential. Lift with $\rho_2E_2$ in both
arguments and average by conjugation exchanging signs; this
preserves positivity, physical commutator and the signed Hadamard
relation. Smooth cut additions propagate smoothly locally by finite
causal shadow. The usual Weyl construction from a real covariance
with $\mu\pm i\sigma/2\geq0$ gives the stated regular state.
This proves Theorem~\ref{thm:generalstate}.

\section{A single reference and modular kernels for all regions}\label{sec:generalmod}
We now choose one state, not a different state for each local algebra.
The state depends on the background, a positive base covariance of
Theorem~\ref{thm:generalstate}, and the construction parameters below.
Its reference is generally mixed, nonstationary and not canonical.
Write $\mathcal V_{\rm r}$ for $\mathcal V$ or its radical reduction,
choosing the corresponding base covariance $\mu_0$.
For $O\subset M$ open, the local phase space is the image in
$\mathcal V_{\rm r}$ of compact Ward sources supported in $O$.
In particular the following theorem applies to every causally convex
diamond and every globally hyperbolic causally convex region, including
unbounded regions.

\begin{theorem}[Bare all-region modular generator]\label{thm:generalmod}
There is a single regular positive Gaussian physical Hadamard reference
such that the restriction to every represented local von Neumann
algebra is faithful and standard in its local cyclic representation.
Every compact smooth Ward source in that region is in the domain of
every power of its one-particle modular Hamiltonian. On each fixed
compact support the graph norms have finite Sobolev bounds. The
one-particle Hamiltonian has a convergent bare odd-power expansion,
and a convergent Abel--Chebyshev reconstruction in the physical Gram
matrices. For relatively compact $O$ the physical logarithmic pairing
has a redundant ordinary distribution kernel on $O\times O$; for
arbitrary $O$ the physical pairing is defined on all compact Ward
tests with fixed-support distribution bounds. No canonical off-Ward
coordinates on a general unbounded region, global
Reeh--Schlieder property, or geometrically preferred reference is
asserted.
\end{theorem}

\subsection{Compact preparation blocks}
Choose auxiliary positive spacetime pairings, lowering the ghost index
by the Lorentzian metric and identifying the resulting bundles
smoothly with the action pairing. Under this identification the Ward
operator is the auxiliary formal adjoint $P_W=K^*$ (distinct from
the Euler operator $P$). At a nonzero real covector $\zeta$ its
adjoint symbol $w\mapsto 2\zeta_{(a}w_{b)}$ is injective, even for
Lorentz-null $\zeta$; the scalar row is zeroth order. Thus $P_W$
is underdetermined elliptic. We describe why the compact support
right inverse of \cite[Theorems 3.8 and 5.1]{delay} applies on
smooth relatively compact collar domains. For its positive
auxiliary weighted estimate compare $\operatorname{sym}\nabla^g$
with the symmetric gradient of an auxiliary Riemannian metric.
After ghost index lowering their difference is order zero
with smooth bounded coefficient near the compact domain boundary.
In the asymptotic weighted Korn estimate this perturbation is
multiplied by the boundary weight $x^2$; shrink the collar and
absorb it. The scalar row increases the positive norm on the
right. Kernel regularity at the boundary follows by elliptic
regularity in the interior and the finite-type Killing connection:
normal transport of the value and antisymmetric first derivative
extends smoothly through the compact boundary collar, also after
tangential differentiation. The scalar row imposes additional
algebraic constraints but does not worsen this estimate. These
verify the asymptotic Poincar\'e and kernel regularity hypotheses
for $P_W$; no absence of joint local Killing fields is assumed.

On a fixed such domain the weighted projected variational inverse
gives a linear right inverse $Q_m$ on sources of a fixed interior
compact support satisfying exact kernel orthogonality, output
flat at the outer boundary. Its ordinary smooth seminorms are
bounded by finite input Sobolev seminorms: use the fixed-domain
weighted inverse and scaled derivative estimates of \cite{delay};
on the fixed source support weighted norms are ordinary-equivalent,
and the boundary multiplier vanishes exponentially, absorbing the
polynomial derivative losses. Extend the output by zero. Choose a
smooth precompact exhaustion with compact cutoffs $\chi_m$ equal
to one on increasing inner sets and supported inside those domains.
For any smooth auxiliary input a, $P_W(\chi_m a)$ is supported
in the prescribed interior compact and satisfies the compatibility
conditions by integration by parts against every homogeneous
$K$-field. Define
\[
 T_m a=\chi_m a-Q_m P_W(\chi_m a).
\]
This is a continuous smooth compact-source Ward preparation, and
$T_m f=f$ for every Ward test supported in the inner set. It is not a
local projection preserving the support of arbitrary non-Ward inputs.

Extend the auxiliary input bundle to a compact double about that
support. Choose a real auxiliary elliptic eigenbasis $e_{mj}$ of
positive order-one frequencies $\lambda_{mj}$ there. Put
$g_{mj}=[T_m e_{mj}]\in\mathcal V_{\rm r}$. For each fixed m their
base covariance norms and smooth bisolution seminorms on any fixed
compact grow polynomially in $\lambda_{mj}$: use compact elliptic
bounds, the fixed-domain continuity just established, and local
smooth Green evolution from a compact source. Since the spectral
count grows polynomially, choose
\[
 \epsilon_{mj}=\theta_m\exp[-\log^2(2+\lambda_{mj})],\qquad
 N=\sum_{m,j}\epsilon_{mj}^2\,
       \sigma(g_{mj},\cdot)\otimes\sigma(g_{mj},\cdot).
\]
Take positive $\theta_m$ small enough that each m-block contributes
at most $2^{-m}$ to the sum of its covariance norms squared with
these weights and to its first m seminorm bounds on a fixed compact
exhaustion of the smooth two-point kernel. This yields a smooth
real bisolution perturbation with
$\sum\epsilon_{mj}^2\|g_{mj}\|_{\mu_0}^2<\infty$. The covariance
$\mu=\mu_0+N$ obeys
$\mu_0\le\mu\le C\mu_0$ by the CCR bound, and keeps positivity,
the same physical charge and Hadamard wavefront. The noise
descends because it uses only $\sigma$-functionals. It can also
be realized by Gaussian Weyl displacements, but no assumption of
full displacement support along central directions is used.

\subsection{Uniform endpoint control on each compact support}
In the real local covariance Hilbert space $Y$ modulo its covariance
null, let $D$ satisfy $\mu(f,Dg)=\sigma(f,g)/2$; then $A=iD$ is a
selfadjoint contraction on $Y_\mathbb{C}$, with conjugation
$J_0 A J_0=-A$. The restricted noise functionals have real Riesz
columns $v_{mj}$, and the bounded column map $V_O$ obeys
$1\pm A\ge V_O V_O^*$ by base CCR positivity. For
$E=1_{[0,\delta]}(1\pm A)$, it follows that
$\|EV_O\|\le \sqrt\delta$. For a fixed compact Ward f, choose
one m-block with $T_m f=f$. Its finite spectral reconstruction
$f_L=[T_m1_{\lambda\le L}f]$ satisfies
$\|f-f_L\|_\mu\le C_m L^{-1}\|f\|_{H^{s_m+1}}$ by continuity
on the compact double. The reconstructed sources need not lie in
the region: their symplectic functionals restrict to $Y$.
The finite combination defining $A f_L$ in that sense uses the
columns of $V_O$ with coefficient norm at most
$C_m\exp(\log^2(2+L))\|f\|_2$.
The CCR bound controls $\|A f-A f_L\|$ by the covariance
approximation error, and $|A|\ge1-\delta$ on $E$. Thus, choosing
$L+2\asymp \exp(\sqrt{|\log\delta|}/2)$, for sufficiently small
$\delta>0$ we obtain
\begin{equation}
 \|Ef\|\le C_m e^{-c_m\sqrt{|\log\delta|}}\|f\|_{H^{s_m+1}}.
                                                               \label{eq:gendpoints}
\end{equation}
On the left only the physical class matters, while on the right
one uses any fixed-support test representative. These constants
are independent of the region as long as it contains the given
support. In particular neither endpoint $\pm1$ is an eigenvalue,
by density of compact labels. Layer-cake integration bounds the
spectral moments of every power of
$k(A)=\log[(1+A)/(1-A)]$ by the same finite test seminorm.

\subsection{Actual Tomita operator and coordinate reconstruction}
The real Weyl embedding j has complex Gram $1+A$. In the
cyclic one-particle space its identification with $Y_\mathbb C$
is $j(f)\mapsto (1+A)^{1/2} f$, onto after completion by
endpoint exclusion. The real image is closed since
$\|j(f)\|^2=\mu(f,f)$ for real f, and its intersection with its
imaginary multiple is zero by injectivity of $(1+A)^{1/2}$.
It is standard in its cyclic complex span. Its Tomita operator is
$J_0[(1-A)/(1+A)]^{1/2}$ on
$\operatorname{ran}(1+A)^{1/2}$, which is exactly the closed
domain by the identity
\[
 \|(1+A)^{-1/2}x\|^2
 =\tfrac12\bigl(\|x\|^2+
    \|[(1-A)/(1+A)]^{1/2}x\|^2\bigr).
\]
Therefore $\Delta_1=(1-A)/(1+A)$ and $-\log\Delta_1=k(A)$.
Bosonic second quantization of the standard subspace gives
the represented modular operator and Hamiltonian
$\Gamma(\Delta_1)$ and $d\Gamma(k(A))$ \cite{bgl,longo}.
In the global Fock representation the local algebra acts on
the Fock space of its cyclic complex span tensored with the
identity on the complementary factor. Thus the restriction is
faithful and its automorphism is the one just computed even
without cyclicity in the entire global representation. Central
represented real directions are in $\ker A$ and have trivial
modular action. This is distinct from a promise of a faithful
representation of an unreduced abstract presymplectic algebra.

With $\delta=n^{-1/2}$ in (\ref{eq:gendpoints}), the endpoint
projection estimate and the decay of $|A|^n$ on its complement
give the absolutely norm-convergent identity
\[
 k(A)f=2\sum_{n\ge0}\frac{A^{2n+1}f}{2n+1}.
\]
The truncation norm is at most
$C_m\exp[-c'_m\sqrt{\log(2+N)}]\|f\|_{H^{s_m+1}}$ after
term N: summing the endpoint bound with weight $1/(2n+1)$
loses at most a polynomial in $\sqrt{\log N}$, absorbed by
decreasing $c'_m$. Alternatively for $0<q<1$ use
\[
 k_q(x)=\log\frac{1+2qx+q^2}{1-2qx+q^2}
  =4\sum_{j\ge1\ {\rm odd}}\frac{q^j}{j}\mathsf T_j(x).
\]
The uniform error at degree J is at most
$4q^{J+1}/((J+1)(1-q))$; $|k_q|\le|k|$ gives the bare limit on
the stated domain. For each fixed polynomial, the finite physical
Gram compression of $A$ converges strongly as its dense
projection basis is exhausted. Hence the ensuing sequential
limits of finite ordinary matrices reconstruct the true generator
in this covariance norm. No fixed finite Gram error uniform
as $q\uparrow1$ is claimed.

The physical form $\mu(a,k(A)b)$ (or the symplectic output
$\sigma(a,k(A)b)$) has fixed-support Sobolev bounds by
(\ref{eq:gendpoints}). For relatively compact $O$, use a single
block containing its closure and project in a sufficiently high
ordinary auxiliary Sobolev norm onto the closure of actual
compact Ward tests in $O$. Compose the two arguments with this
projection and apply the Schwartz kernel theorem. This gives
an ordinary redundant bundle distribution and its restricted
physical form is independent of the redundant extension.
One must not apply a support-enlarging $T_m$ to an arbitrary
non-Ward test and treat its image as a local vector in $O$.
For unbounded $O$ the same argument gives fixed-support
representatives and the physical form on all compact Ward tests
with the actual single $A_O$. This proves
Theorem~\ref{thm:generalmod}. The round example below has
canonical material output coordinates and the more explicit
Gram kernel in those coordinates.



\section{Worked case: the closed active four-scalar background}\label{sec:round}

The positivity of a physical Hadamard covariance and the determination of a local modular operator are distinct problems. The compact-Cauchy construction of positive covariances for critical Einstein--scalar multiplets in \cite{p09} provides a state on a coupled physical quotient. The flat Weyl calculation in \cite{p21} instead starts from a specific conformal vacuum and a curvature-only algebra. We ask how much information local curvature carries on an actual active background of the first construction, and what can be computed about the modular operator of the resulting \emph{full} local algebra. In this worked case we use the original compact-cut reference of \cite{p09}, independently of the more general existence theorem above; it is not used as a modular or local-generation theorem. General Peierls and quasifree modular machinery is available in \cite{hack,khavpeierls,longo}; the closed-background local rigidity, preparations, and compact-test endpoint estimates below are the additional inputs.

In units \(8\pi G=1\), take the action and exact closed FLRW solution
\begin{equation}
 S=\int\left(\frac{R}{2}-\frac12\sum_{A=1}^4|d\Phi^A|_g^2\right)\vol_g,\qquad
 g=-dt^2+a(t)^2\gamma,\qquad \Phi^A=f(t)n^A .
 \label{eq:bg}
\end{equation}
Here \(\gamma\) is the round unit metric on \(S^3\) and \(n:S^3\hookrightarrow\mathbb R^4\) is the unit embedding. Thus \(n\cdot n=1\), \(dn\cdot dn=\gamma\), and \(\Delta_\gamma n=-3n\). Put \(H=\dot a/a,\ b=\dot f/f,\ z=a^{-2},\ F=f^2\). A dot denotes proper-time differentiation. The exact background equations are
\begin{align}
 3(H^2+z)&=\tfrac12 F(b^2+3z),&
 \dot H&=z-\tfrac12 F(b^2+z),\nonumber\\
 \dot b&=-3Hb-3z-b^2,&
 \dot z&=-2Hz,& \dot F&=2bF. \label{eq:ode}
\end{align}
They follow from \(\mathrm{Ric}_{\mu\nu}=\partial_\mu\Phi\cdot\partial_\nu\Phi\) and \(\Box_g\Phi=0\). The constraint in \eqref{eq:ode} propagates. We use the actual example of \cite{p09}, with
\begin{equation}
 (a,f,H,\dot f)(0)=(1,1,1,3).
 \label{eq:seed}
\end{equation}
The regular ODE gives a smooth globally hyperbolic slab \(I\times S^3\). We shorten \(I\) below, retaining \(t=0\). There is no spectator field. The potential is identically zero, the target metric is positive and fixed, and \(S^3\) is compact Cauchy, so the hypotheses of \cite{p09} hold. The material map \(\Phi=f(t)n\) is a diffeomorphism onto an annulus after choosing \(f,\dot f>0\); the metric--scalar Hessian is nonzero, for example \(\nabla_t\nabla_i\Phi^A=(\dot f-Hf)\partial_i n^A\) at the cut.

We use the common metric--scalar pairing and gauge normalization of \cite{p09}. Denote its formally self-adjoint Euler operator by \(P\), its infinitesimal diffeomorphism generator by \(K\), and its formal gauge adjoint by \(K^\dagger\). The normalizations can be translated back into raw metric and scalar variations by the invertible trace-reversal map of that paper. The identities used here are \(PK=K^\dagger P=0\), and the normally hyperbolic pair
\[
 D_2=P+KK^\dagger,\qquad D_1=K^\dagger K,\qquad D_2K=KD_1.
\]
Write \(E_i=G_{i,\mathrm{ret}}-G_{i,\mathrm{adv}}\), choosing the corresponding common Peierls convention. A compact Ward source is a smooth compact source \(j\) with \(K^\dagger j=0\). Its physical solution is \([E_2j]\), modulo linearized diffeomorphisms, and \(\sigma\) is the induced physical symplectic form \cite{p09,hack,khavpeierls}. In each region the local algebra is generated by the \emph{ambient} physical classes of its compact Ward sources. We never identify independently defined regional quotients by assuming an extra support property.

Let \(\mathcal A(O)\) denote this local von Neumann algebra in the specified representation, and \(\mathcal W(O)\) the subalgebra generated by the compact linearized Weyl smears. Since the background Weyl tensor vanishes, \(\delta C\) is gauge invariant. All the results below concern the linearized algebra on this exact background. In particular a modular generator below is not the nonlinear material-time Hamiltonian.

\section{One scalar completes the local curvature algebra}

Let \(r_\Phi=|\Phi|\) and define in material labels
\begin{equation}
 \mathcal N(y)=
 \left[g^{\mu\nu}\partial_\mu r_\Phi\,\partial_\nu r_\Phi\right]
       \bigl(\Phi^{-1}(y)\bigr).
 \label{eq:N}
\end{equation}
On a compact label support the inverse need only be defined on the branch near the background, with a buffered collar. Its first variation is a finite differential, gauge-invariant readout; support in the following statements means support in the background preimage of the compact label test. The branch dependence is the usual one for generalized local gravitational observables \cite{khavlocal}.

\begin{theorem}[Local completion and strictness]\label{thm:completion}
After shortening \(I\), on every open \(O\subset I\times S^3\), in any distributional quasifree physical representation, the compact tests of \(\delta C\) and \(\delta\mathcal N\) generate \(\mathcal A(O)\). If \(O\) is nonempty, then \(\mathcal W(O)\subsetneq\mathcal A(O)\). The equality is one of closed real one-particle spaces and hence of represented von Neumann algebras; no algebraic differential inverse of compact field tests is asserted.
\end{theorem}

\subsection{Local rigidity including distributional solutions}

We first show that an on-shell distributional perturbation with both readouts zero in a chart is pure gauge there. By local conformal flatness, vanishing linearized Weyl implies a conformal representative \(h=2\psi g\) modulo a Lie derivative. For completeness, in a conformally flat chart reduce to the flat linearized Weyl equation. The linearized Bianchi identity gives zero curl of the linearized Schouten tensor. Two local Poincar\'e integrations express it as a Hessian. Subtract the conformal perturbation with that Schouten tensor; the residual linearized Riemann tensor is zero and two flat compatibility integrations give a symmetrized derivative. These integrations apply to distributions on smaller star-shaped patches by the distributional de Rham homotopy. The conformal covariance of linearized Weyl at a conformally flat background brings the splitting back to \(g\). No global choice of the conformal gauge is needed.

Write the scalar variation as \(\delta\Phi=A n+dn(V)\), with \(V\) a spatial vector for the round metric, and write \(D\) for the round connection. In this representative the trace-reversed Einstein equation gives
\begin{align}
 -3\ddot\psi-3H\dot\psi+z\Delta_\gamma\psi
    &=2\dot f\,\dot A,\nonumber\\
 -2D_i(\dot\psi-H\psi)
    &=\dot f(D_iA-V_i)+f\dot V_i,\nonumber\\
 -2D_iD_j\psi+
  a^2(\ddot\psi+5H\dot\psi-z\Delta_\gamma\psi)\gamma_{ij}
    &=2f(A\gamma_{ij}+D_{(i}V_{j)}).
 \label{eq:weylzero}
\end{align}
Indeed \(\delta\mathrm{Ric}_{\mu\nu}=-2\nabla_\mu\nabla_\nu\psi-g_{\mu\nu}\Box_g\psi\), and differentiation of \(dn(V)\) in a spatial direction produces normal component \(-V_i n\). Rank four also means that in this setting the linearized scalar equations follow from the linearized Bianchi identity once all Einstein equations hold. At fixed material labels the compensated scalar reading is
\begin{equation}
 e=\frac{\delta\mathcal N}{2\dot f^2}
   =\psi-\partial_t(A/\dot f).
 \label{eq:e}
\end{equation}
Indeed at the original spacetime point the variation is \(2\dot f^2\psi-2\dot f\dot A\), and the inverse label variation adds \(2\ddot f A\).

The tracefree spatial equation in \eqref{eq:weylzero} gives \(V=-D\psi/f+V_{\rm CK}\). The local spatial conformal Killing fields on \(S^3\) are the restrictions of six rotations and four gradients \(Dn^A\); this follows by the finite-type flat conformal Killing equation in a conformal chart. With distributional time dependence, test against compact temporal functions and extract the same finitely many spatial coefficients by compact spatial dual tests. Taking the spatial curl of the mixed equation shows that the rotation part satisfies \(\dot V_{\rm rot}=bV_{\rm rot}\) and is pure gauge. There is no kernel of this curl on the round local rotation space: a closed Killing one-form would be parallel, and the round curvature then forces zero.

Set \(\mathcal E=\operatorname{span}(1,n^1,\dots,n^4)\) restricted to the patch. The mixed equation now gives \(A\equiv(-\dot\psi+c\psi)/\dot f\) modulo \(\mathcal E\), with \(c=2(H-b)\). All remaining conformal Killing contributions to the scalar trace equations lie in \(\mathcal E\). Put \(d=\ddot f/\dot f=-3H-3z/b\). Substitution in the normal and spatial trace equations gives, modulo \(\mathcal E\),
\begin{align*}
 \ddot\psi+(5H+2z/b)\dot\psi-z\Delta_\gamma\psi-2zc\psi/b&\equiv0,\\
 \ddot\psi+(3H+2c+2d)\dot\psi-z\Delta_\gamma\psi
                       +(2\dot c-2dc)\psi&\equiv0.
\end{align*}
In particular, if
\begin{align*}
 D_0&=-4(H+b+2z/b),\\
 E_0&=2\dot c-2dc+2zc/b,\\
 r_0&=-E_0/D_0,\\
 k_0&=-\bigl[\dot r_0+r_0^2+(5H+2z/b)r_0-2zc/b\bigr]/z ,
\end{align*}
then \(\dot\psi\equiv r_0\psi\), \((\Delta_\gamma+k_0)\psi\equiv0\). Differentiate the latter equation; the terms containing \(\dot\psi\) cancel modulo \(\mathcal E\), leaving \(\dot k_0\psi\equiv0\). All these operations are valid on distributions; multiplication by background time functions and the spatial Laplacian preserve the removed subspace. At \eqref{eq:seed}, direct differentiation by \eqref{eq:ode} gives
\[
 D_0=-56/3,\quad E_0=100/3,\quad
 r_0=25/14,\quad \dot r_0=-3191/196,\quad
 k_0=15/49,\quad \dot k_0=-20835/686 .
\]
Thus on a smaller common slab \(\dot k_0\ne0\) and \(\psi,A\) are in \(\mathcal E\). This argument is local: no harmonic boundary condition or unique continuation of an arbitrary elliptic solution is used at a spatial rim.

The resulting low-mode coefficients solve smooth ordinary systems and hence are smooth, also for distributional input. We give the rank check that \eqref{eq:e} detects all of them modulo gauge. In either the homogeneous or dipole block, set \(p=\psi\), \(v=\dot\psi\), \(x=A/f\); for a dipole \(V=B Dn^A\) set \(y=B/f\). The equations restricted to these coefficients are
\begin{align*}
 \dot p&=v,\quad
 \dot v=-5Hv+2Fzx\quad\text{(homogeneous)},\\
 \dot v&=-5Hv-5zp+2Fz(x-y),\quad
 \dot y=-2(v-Hp)/F-bx\quad\text{(dipole)},\\
 \dot x&=(-3\dot v-3Hv-kzp)/(2Fb)-bx
     \quad(k=0,\ 3 \text{ respectively}).
\end{align*}
At the cut residual spacetime conformal Killing freedom sets \(p=0\) in the homogeneous harmonic, leaving \((v,x)\). For the dipole it sets \(p=v=0\), leaving \((x,y)\). To see the latter rank, in background conformal time the dipole generator has coefficients \((U_0,V_0)\) with \(U_0'=-V_0,\ V_0'=U_0\); its metric conformal factor and proper-time derivative at the cut are \(U_0-V_0,-4U_0-V_0\), independent. The homogeneous conformal time translation has conformal factor \(aH\ne0\). Differentiating \(e=p-\partial_t(x/b)\) with the displayed equations gives at the public cut
\[
\begin{array}{c|cc}
 & e & \dot e\\ \hline
 \text{homogeneous}&-(2v+3x)/3&-(v+x)/3\\
 \text{dipole}&-(3x+y)/3&(2x-11y)/3 .
\end{array}
\]
The determinants are \(-1/9\) and \(35/9\). They stay nonzero on a sufficiently small common time slab by smooth ODE propagation on the two-dimensional blocks. Angular independence on any open chart then shows that \(e=0\) throughout the chart forces all its low physical coefficients to vanish. In a distributional chart the coefficients are smooth because a smooth ODE fundamental matrix reduces their distributional derivative to zero. We have proved local gauge triviality when both readings vanish. The gauge vector can be taken to be the unique \((d\Phi)^{-1}\delta\Phi\) wherever there is gauge triviality, so these identifications are compatible on overlaps.

\subsection{From rigidity to equality of closed real local spaces}

Write \(\mathfrak k[j]\) for the one-particle vector of a source class in the chosen quasifree representation. On raw perturbations let \(B_0(h,\delta\Phi)=(d\Phi)^{-1}\delta\Phi\), with the corresponding normalization if using the rescaled gauge generator, and \(Q_0=1-KB_0\). Then \(Q_0^*\) takes arbitrary compact tests to Ward sources, preserves support, and is the identity on Ward sources. Let \(\zeta\) be real-orthogonal to the closed real span of the two readout families in \(O\). Define the distribution
\[
 u_\zeta(j)=\operatorname{Re}\langle\zeta,\mathfrak k[Q_0^*j]\rangle .
\]
It is a distribution by continuity of the two-point function on compact tests. It solves \(Pu_\zeta=0\) because \(Pq\) is Ward closed and represents zero. The assumed orthogonality says the two differential readouts vanish on \(u_\zeta\) throughout \(O\). Local rigidity therefore makes \(u_\zeta\) pure gauge there, and it annihilates every Ward source supported in \(O\) by integration by parts. A cutoff of the gauge vector equal to it near the source makes this assertion valid without any global support assumption on the potential gauge vector. The two closed real spans have the same real orthogonal complement and hence are equal. Strong continuity of the Weyl representation gives equality of the von Neumann algebras.

To prove strictness take a nonzero exact target translation \(\Phi\mapsto\Phi+\epsilon v_0\). It leaves the metric, and therefore Weyl, unchanged. In the low-mode variables \(x=y=(v_0\cdot n)/f\); at \eqref{eq:seed} the dipole reading is \(-4(v_0\cdot n)/3\). Its nonzero coefficient persists on a thinner time slab, so it changes a compact radial reading in any nonempty open set there. The translation tangent is a smooth global physical solution and is represented by a compact Ward source, since its Cauchy surface is compact and the gauge-compatible time-slice construction applies \cite{p09,hackschenkel}. Its Weyl implementer commutes with every curvature unitary in the particular region but not with a suitably scaled local radial unitary, by the nonzero physical phase of this reading. This proves strictness in every regular quasifree realization and completes Theorem~\ref{thm:completion}.

\section{Local preparations and information lost by curvature restriction}

\begin{theorem}[Properly localized information separation]\label{thm:local}
There are arbitrarily small nested proper concentric diamonds \(O_1\Subset O_2\) in the above slab, and a real two-parameter family \(w_m\) of physical displacements implemented by Weyl unitaries in \(\mathcal A(O_2)\). They fix \(\mathcal W(O_1)\) pointwise under conjugation. There are two fixed compact coupled readouts in \(O_1\) whose mean shifts are \(Rm\) for a nonsingular real matrix \(R\). In any centered regular quasifree reference with finite readout covariance \(B\), the corresponding coherent pair has zero relative entropy on \(\mathcal W(O_1)\) and coupled relative entropy at least \(\tfrac12m^TR^TB^{-1}Rm>0\) when \(m\ne0\). The preparers may be chosen with nonzero mutual symplectic product.
\end{theorem}

\subsection{Exact compact data and actual local representatives}

In material coordinates at \(t=0\) fix \(\Phi=f(0)n\); the conjugate variables are the induced metric \(q\), \(\pi^{ij}=w(K^{ij}-Kq^{ij})/2\), and \(P_A=w U_A\), where \(w=\sqrt{\det q}\), \(K=\mathrm{tr}_q K_{ij}\) is the expansion with \(K_{ij}=+\mathcal L_n q_{ij}/2\), and \(U_A\) is the scalar proper normal derivative. The constraints are
\begin{align}
 U_A\partial_i\Phi^A&=2q_{ij}D_k\pi^{jk}/w,\\
 |U|^2&=R(q)+K^2-K_{ij}K^{ij}-q^{ij}\partial_i\Phi\cdot\partial_j\Phi .
 \label{eq:constraints}
\end{align}
The first determines the three tangential components of \(U\) by the material Gram inverse; the second fixes the positive radial component by a square root after subtracting the tangential square. The branch is open at the seed because the squared radial momentum is \(\dot f^2>0\). This solve is local in the initial data, including their displayed derivatives. The rank of the four material derivatives means that prescribing \(\delta\Phi=0\) on the cut fixes all four infinitesimal gauge deformations of the Cauchy data. The restriction of the ADM symplectic form is therefore exactly the canonical form on \((\delta q,\delta\pi)\); the scalar term is \(\delta P_A\wedge\delta\Phi^A\) and disappears on this slice. This is a reduction, not an assumption of a separated test-graviton symplectic form.

Use conformal time \(\eta=\int_0^t a^{-1}ds\) and small geodesic balls about a center, of radii \(0<r<\rho<b<R_0\). Take \(O_1=\{d_\gamma+|\eta|<r\}\), \(O_2=\{d_\gamma+|\eta|<R_0\}\), with closures in the slab and radii smaller than the injectivity scale. They are proper causally convex diamonds. Consider the exact homogeneous data varying \(a,H\) by tangent \(m=(m_a,m_H)\) at fixed \(f(0)\); their positive \(\dot f\) is solved by the Friedmann constraint. Call their metric/momentum data \(q_{\rm h}(s),\pi_{\rm h}(s)\) along \(sm\). Choose a real smooth bump \(\beta=1\) near the closed \(\rho\)-ball, supported in the \(b\)-ball. Cut off only the metric data:
\[
 q_s=\bar q+\beta(q_{\rm h}(s)-\bar q),\qquad
 \pi_s=\bar\pi+\beta(\pi_{\rm h}(s)-\bar\pi).
\]
Then apply the \emph{exact} scalar momentum solve \eqref{eq:constraints}. The full constrained data agree with the homogeneous data on the \(\rho\)-ball and with the seed outside the \(b\)-ball. In particular differentiating \(\beta\) has not been ignored in the scalar momenta. Local Einstein--scalar Cauchy existence and uniqueness imply that after differentiation the solutions agree with the homogeneous variation in its domain of dependence, and are supported, in a suitable gauge representative, in \(J(\overline B_b)\). One can see both facts by using a common background-matched wave-coordinate gauge: choose the same initial lapse and shift values, use the local data to fix their first derivatives by the gauge condition, and apply causal uniqueness to the resulting linear hyperbolic system. The constraints give the subsidiary initial conditions. Converting this supported representative to \(K^\dagger w_m=0\) is done by solving the ghost equation \(D_1\xi=-K^\dagger w_m^{\rm old}\) with vanishing initial data; it preserves the same causal support. The entire homogeneous family is conformally flat, so the \emph{full} linearized Weyl tensor of \(w_m\) vanishes on \(O_1\).

Take a temporal switch \(\chi\) from zero to one with derivative supported in \(|\eta|\le\delta\), where \(b+2\delta<R_0\). Then \(j_m=P(\chi w_m)\) is a compact Ward source supported in \(O_2\). Indeed the support lies in the switch band inside \(J(\overline B_b)\), and \(K^\dagger P=0\). The usual retarded and advanced inverse identities on \(\chi w_m,(\chi-1)w_m\) give
\begin{equation}
 E_2D_2(\chi w_m)=w_m,\qquad
 E_2j_m=w_m-KE_1K^\dagger(\chi w_m).
 \label{eq:switch}
\end{equation}
The last input is compact since \(K^\dagger w_m=0\). Thus the Weyl operator really belongs to \(\mathcal A(O_2)\) and implements the indicated displacement, rather than merely a globally defined shift of expectations.

\subsection{Two actual readouts and finite entropy bounds}

Let \(n_\Phi=-\nabla r_\Phi/\sqrt{-g(\nabla r_\Phi,\nabla r_\Phi)}\) be the material future normal. The material observables corresponding to the homogeneous radius and expansion are
\[
 \mathcal A_{\rm r}=
 \left[\frac{3r_\Phi^2}
 {(g^{\mu\nu}+n_\Phi^\mu n_\Phi^\nu)\partial_\mu\Phi\cdot\partial_\nu\Phi}
 \right]^{1/2},\qquad
 \mathcal H=\nabla_\mu n_\Phi^\mu/3 .
\]
Evaluate at fixed material label as in \eqref{eq:N}. On the seed these are \(a,H\); their fixed-clock variations subtract \((\dot a/\dot f)\delta r_\Phi\), \((\dot H/\dot f)\delta r_\Phi\) respectively from the coordinate variations. The local perturbations above have values \(m_a,m_H\) at the cut. A fixed normalized positive compact smooth density close enough to that cut and supported in \(O_1\) gives an integrated response \(R\) close to the identity. It may be fixed with \(R\) nonsingular. There is active matter, not a spectator perturbation: at the public cut, for example, the homogeneous direction \(\delta a=\delta f=0,\ \delta H=1/2\) has \(\delta\dot f=1\) by \eqref{eq:ode}. The same sign branch throughout the preparation ensures that this actual material momentum, when measured in the inner region, is the homogeneous response.

The conjugation by \(\mathrm W(j_m)\) fixes every curvature unitary in the inner algebra, by its zero physical phase there, hence fixes \(\mathcal W(O_1)\) pointwise. For the two basic preparers the material canonical formula gives
\[
 \sigma(w_a,w_H)=-\frac{6}{a(0)}
            \int_{S^3}\beta^2\,\vol_{\bar q}\ne0
\]
in the convention \(\sigma=\int(\delta q_1\delta\pi_2-\delta q_2\delta\pi_1)\). In fact \(\delta_a q=2\beta\bar q/a\), \(\delta_H q=0\), and \(\delta_H\pi=-\beta\bar w\bar q^{-1}\), which suffice for this pairing. Thus the relative commutant \(\mathcal W(O_1)'\cap\mathcal A(O_2)\) contains a noncommuting pair (rescale to avoid a trivial Weyl phase), even though their conjugations commute as automorphisms.

For a quasifree reference the two coupled readouts have a finite symmetric covariance \(B\), positive definite by CCR uncertainty since any nonzero readout combination pairs nontrivially with some \(j_m\). Measuring the single linear combination with coefficient \(u\) gives Gaussians with common variance \(u^TBu\) and mean shift \(u^TRm\). The classical relative entropy and its monotonicity under measurement give
\[
 S_{\mathcal A(O_1)}(\omega_m\mid\omega)\ge
 \sup_u\frac{(u^TRm)^2}{2u^TBu}
 =\frac12 m^TR^TB^{-1}Rm .
\]
The readouts need not commute because only a single combination is measured in each bound. The reference and local unitary-shifted state are normal in the same representation. This proves Theorem~\ref{thm:local}; the lower bound does not assert that the full entropy of an arbitrary coherent pair is finite.

For a finite two-sided comparison, use dual readout coordinates in which \(R=1\), write \(\beta_x=\operatorname{Ad}\mathrm W(j_x)\) with the sign giving the chosen mean shift, and put \(\omega_d=\int(\omega\circ\beta_x)\,dG_{d,C}(x)\) for a real two-dimensional Gaussian of covariance \(C>0\). These are normal quasifree states: the implemented predual orbit is norm continuous, and its Weyl characteristic function adds exactly the smooth finite-rank physical covariance of the preparations. Their Weyl marginals agree. The UCP map \(A\mapsto[x\mapsto\omega(\beta_x(A))]\) to Lebesgue \(L^\infty\) is normal by the predual integral. The corresponding upper data-processing bound, and the same single-readout lower bound with variance \(B+C\), give
\begin{equation}
 \frac12 d^T(B+C)^{-1}d
 \ \le S_{\mathcal A(O_1)}(\omega_d\mid\omega_0)
 \ \le \frac12 d^TC^{-1}d,\qquad
 S_{\mathcal W(O_1)}(\omega_d\mid\omega_0)=0 .
 \label{eq:finite-entropy}
\end{equation}
Both orientations are bounded in the same way. Thus the restriction loses finite strictly positive information when \(d\ne0\); with \(C=\tau C_0\), the entropy times \(\tau\) tends to \(d^TC_0^{-1}d/2\). For a faithful reference, these mixtures are faithful, and Petz's sufficiency criterion \cite[Theorem 4]{petz} also implies that their full relative Connes cocycle is not contained in the inner Weyl algebra, whereas the restricted cocycle is the identity. That statement alone says nothing about preservation by an arbitrary absolute modular group.

\section{A faithful Hadamard reference and the bare diamond modular kernel}

We give a state-specific construction on the same \(\mathcal A(O)\), with \(O\) a proper cut diamond. It uses all active physical data, which Theorem~\ref{thm:completion} equivalently generates by Weyl and \(\delta\mathcal N\). First identify its real local smooth labels with the canonical smooth cut pairs supported compactly in the cap \(B\): compact sources have such support by causal propagation and the support-preserving material gauge. Conversely any such pair lifts by \eqref{eq:constraints} to constraint-compatible data; the switch argument \eqref{eq:switch}, with a band inside \(O\) using the support buffer, produces a compact Ward source in \(O\) representing it. Neither implication uses a spatial TT projection.

\subsection{Specified reference and spectral endpoint bound}

Start with one physical Cauchy covariance of \cite{p09}, fixing its collar, frequency-projection smoothing choices, and finite-dimensional Fredholm correction. For example its real covariance is the realified mean of the two forms \(\lambda_\Sigma^\pm=\pm T^*Q c_2^\pm T+\nu_\pm\) of \cite{p09}. These are state inputs and are not selected by the metric. Average the two-point \emph{covariance} over the compact diagonal \(SO(4)\) acting by round spatial rotations and matching target rotations at the seed, and take the quasifree state of the resulting covariance \(\mu_{\rm b}\). The commutator and positivity persist. The Hadamard wavefront inclusion persists under smooth joint pullback by the group action and compact pushforward over the group: every surviving spacetime wavefront pair is a rotated future-null pair in the original relation. This is a distributional averaging argument, not an assumption of uniformity from pointwise smooth remainders. Smooth physical differences remain smooth under the same averaging. The covariance is continuous in a finite Sobolev norm of the global smooth canonical Cauchy pair, because the cut lift to compact Ward tests with a fixed switch is continuous in the smooth topologies, and the original covariance is distributional.

Let \(C_{\rm ang}\) be the positive group Casimir on the global material Cauchy labels. It is elliptic on the compact cut, since the rotational Killing generators span each tangent space, and has finite-dimensional real modes of polynomial spectral count. Let \(J_{\rm ang}=(1+C_{\rm ang})^{1/2}\) with projections \(\Pi_\lambda\) for its levels \(\lambda\ge1\). The canonical form and invariant covariance are block diagonal in those levels. Fix a number \(\theta>0\) and choose the reference
\begin{equation}
 \mu=\mu_{\rm b}+N,\qquad
 N(f,g)=\theta\,\mu_{\rm b}(f,e^{-J_{\rm ang}}g).
 \label{eq:noise}
\end{equation}
In auxiliary ordinary canonical \(L^2\) pair coordinates, each finite block has the canonical invertible symplectic matrix. Thus \(N\) is actual centered Gaussian displacement noise: on each block take the positive Gaussian covariance obtained by conjugating its noise matrix with the inverse transpose and inverse of the canonical symplectic matrix. The covariance coefficients grow at most polynomially in the unnoised basis, and the exponential factor makes the added Cauchy kernel smooth in all derivative orders. Exact physical Cauchy propagation preserves this smooth bisolution addition. This verifies positivity, physical CCR and Hadamard regularity on the same gauge quotient. The Gaussian in the global covariance completion has full support (positive noise on each finite block and independent summably small tails). Consequently the normal Gaussian average in the original cyclic representation is faithful: if a positive operator has zero averaged expectation its square root annihilates the Weyl orbit by full support and strong continuity, and that orbit has dense linear span. We may therefore use the faithful cyclic \emph{local} representation of \eqref{eq:noise}. No unproved local Reeh--Schlieder assertion for the original state is needed.

Let \(L\) be the real local covariance completion, identified with the closed real subspace of the global \(\mu\)-completion generated by local labels. Complexify the real completion, with inner products linear in the second argument. The skew operator \(D\) defined by \(\mu(f,Dg)=\sigma(f,g)/2\) has norm at most one by CCR positivity. Put \(A=iD\), a self-adjoint contraction. In the cyclic one-particle representation the embedding has Gram form \(1+A\); canonical conjugation sends \(A\) to \(-A\). The Tomita factorization therefore gives
\begin{equation}
 \Delta_1=\frac{1-A}{1+A},\qquad
 \mathcal K_1=-\log\Delta_1=k(A),\qquad
 k(x)=\log\frac{1+x}{1-x}.
 \label{eq:polarizer}
\end{equation}
More explicitly identify the physical complex one-particle space with the ordinary complexification by the map \(f\mapsto(1+A)^{1/2}f\); Tomita is \((1+A)^{1/2}J_0(1+A)^{-1/2}\) with canonical conjugation \(J_0\). Faithful local standardness excludes eigenvalues at \(\pm1\): a nonzero such eigenvector and its conjugate would give nonzero imaginary multiples in the local real subspace. The endpoint estimate below also proves their absence. No endpoint spectral gap is assumed. The bosonic Fock modular Hamiltonian is the second quantization \(d\Gamma(\mathcal K_1)\), including zero on the cyclic vacuum \cite{araki,longo,longofinite}; our convention for the automorphism is \(\Delta^{it}\).

\begin{proposition}[Bare log domain]\label{prop:domain}
For the reference \eqref{eq:noise}, every compact smooth material Cauchy label in the cap lies in \(\mathrm{Dom}(\mathcal K_1^d)\) for each positive integer \(d\). The resulting graph norm is bounded by a finite Sobolev seminorm of that label on the cut.
\end{proposition}
\begin{proof}
On the global \(\mu\)-space the noise is represented by \(B=\theta e^{-J_{\rm ang}}/(1+\theta e^{-J_{\rm ang}})\). On the local complexified space set \(T_\pm=1\pm A\). Positivity of \(\mu_{\rm b}\pm i\sigma/2\) gives \(\langle x,T_\pm x\rangle_\mu\ge\langle x,Bx\rangle_\mu\). This is a form inequality with global \(B\) restricted to local \(x\), not a commutation of an angular cutoff and localization. If \(x\) is a unit vector in the spectral band \(T_\pm\le e^{-s}\), then \(\|\Pi_{\le r}x\|_\mu\le C_\theta e^{(r-s)/2}\), by the minimum weight of \(B\) on \(\lambda\le r\). Thus for a smooth local \(f\), also viewed as its smooth global cut pair,
\[
 \|E_{T_\pm}([0,e^{-s}])f\|_\mu
 \le C_\theta e^{(r-s)/2}\|f\|_\mu+\|\Pi_{>r}f\|_\mu .
\]
No support property of \(\Pi_{\le r} f\) is needed in this pairing. Set \(r=s/2\), and use
\(\|\Pi_{>r}f\|_\mu\le r^{-m}\|J_{\rm ang}^m f\|_\mu\) for arbitrary \(m\). The right-hand side is controlled by a finite global cut Sobolev seminorm of \(f\) by ellipticity and covariance continuity. The spectral layer-cake integral for the \(2d\)-th log moment is bounded by integrating \(s^{2d-1}\) against the square of this tail; it is finite with \(m=d+1\). Applying the two endpoint estimates to \eqref{eq:polarizer} proves the claim.
\end{proof}

\subsection{An ordinary coordinate distribution, with a convergent series}

Fix a nested dense real sequence of compact Cauchy labels (or use the two compact readout species of Theorem~\ref{thm:completion} and their physical causal Cauchy data). Gram--Schmidt in \(\mu\) gives \(e_j\), skipping zero residuals, and the finite polarizer matrices and coordinate functionals are
\[
 A_N=\bigl(i\sigma(e_i,e_j)/2\bigr)_{i,j\le N},\qquad
 \phi_j(f)=\mu(e_j,f).
\]
Here the Gram data use the coupled \(D_2\) Cauchy propagator, physical charge, and \eqref{eq:noise}; they are not flat Weyl symbols. Use the notation \(A_N\) also for the compression extended by zero on its orthogonal complement. Then \(A_N\to A\) strongly and \(\|A_N\|\le1\). The real local canonical \emph{coordinate output} \(e_i^{\rm out}\) is the distribution
\(e_i^{\rm out}(\psi)=\sigma(y_\psi,e_i)\), where \(y_\psi\) is the ordinary canonical symplectic conjugate of a compact coordinate output test \(\psi\), with the fixed cut density. The same expression tests output for any element in the local covariance completion, by the global CCR bound.

For real \(0<q<1\) let \(\mathsf T_n\) be the Chebyshev polynomials. The identity
\begin{equation}
 k_q(x)=\log\frac{1+2qx+q^2}{1-2qx+q^2}
       =4\sum_{\substack{n\ge1\\ n\ {\rm odd}}}
          \frac{q^n}{n}\,\mathsf T_n(x)
 \label{eq:abel}
\end{equation}
follows by the Fourier-log series at \(x=\cos\vartheta\). It is uniformly convergent on \([-1,1]\) for fixed \(q\); the tail after \(J\) is at most \(4q^{J+1}/((J+1)(1-q))\). The polynomial compression converges strongly at fixed \(q,J\). On \((-1,1)\), \(k_q\to k\) and \(|k_q|\le |k|\), since \(2q/(1+q^2)\le1\). Proposition~\ref{prop:domain} therefore gives convergence to the \emph{bare} \(k(A)f\) in covariance norm on compact test inputs. An ordinary coordinate distribution kernel, interpreted complex bilinearly or with the corresponding conjugation in a Hilbert-space pairing, is
\begin{equation}
 \boxed{
 \mathcal K_{\rm coord}(x,y)=
 \lim_{q\uparrow1}\lim_{J\to\infty}\lim_{N\to\infty}
  \sum_{i,j\le N}
  \left[4\sum_{\substack{n\le J\\n\ {\rm odd}}}
        \frac{q^n}{n}\mathsf T_n(A_N)\right]_{ij}
  e_i^{\rm out}(x)\otimes\phi_j(y) .}
 \label{eq:kernel}
\end{equation}
The odd polynomial vanishes on the zero complement of the finite compression. Testing output by \(\sigma(y_\psi,\cdot)\) gives precisely the usual material Cauchy-coordinate pairing, rather than only a covariance-dual expression. A uniform finite Sobolev estimate on the compact test pairs follows from the log graph bound and the CCR estimate on \(y_\psi\). It supplies the asserted distributional limit: at fixed polynomial stage the operator norms are bounded; for fixed \(q\) the series is uniform; and at the final stage the graph estimate dominates the matrix elements. Tensor test density and the Schwartz kernel theorem give the corresponding distribution on the coordinate product. Multiplication by the fixed factor dictated by \(\Delta^{it}\) yields the real infinitesimal modular action. No local prescription for a product of singular kernels is required in \eqref{eq:kernel}.

Equation~\eqref{eq:kernel} is a constructive convergent representation of the coordinate distribution for the \emph{specified} reference, not an evaluated elementary spacetime Green function or a universal local modular energy density. The state and the fixed region are essential inputs. The compact-test endpoint claim does not say that these tests are a core for the maximal logarithmic domain.

\section{Discussion of the round worked case}

The exact local generation theorem identifies a minimal addition to Weyl \emph{on this active stock}: one material scalar suffices for the represented net, while Weyl alone fails even in arbitrarily small regions. The localized preparations show that the missing information is physically usable within a proper outer region, not just a global homogeneous mode or a spectator. Equal Weyl marginals can conceal finite full relative entropy. The modular operator reconstructed here is consequently that of the coupled algebra, not the flat Weyl modular generator of \cite{p21}; no transfer of that numerical formula is asserted. The covariance construction, local proof, and coordinate reconstruction all concern linearized physical observables. They make no claim of interacting constraint quantization, anomaly cancellation, a nonlinear modular net, or ultraviolet completion.

\paragraph{Provenance and review.}
This manuscript joins two separate archived research inputs on a common explicitly verified Einstein--scalar background. It was prepared with AI assistance and internally cross-checked against the archived source calculations; no external peer review or priority certification is asserted.

\begin{thebibliography}{99}
\bibitem{bgl} R. Brunetti, D. Guido and R. Longo, \emph{Modular localization and Wigner particles}, Rev. Math. Phys. 14 (2002), \href{https://arxiv.org/abs/math-ph/0203021}{arXiv:math-ph/0203021}.
\bibitem{costabel} M. Costabel and A. McIntosh, \emph{On Bogovski\v\i{} and regularized Poincar\'e integral operators for de Rham complexes on Lipschitz domains}, Math. Z. 265 (2010), \href{https://arxiv.org/abs/0808.2614}{arXiv:0808.2614}.
\bibitem{delay} E. Delay, \emph{Smooth compactly supported solutions of some underdetermined elliptic PDE, with gluing applications}, Comm. Partial Differential Equations 37 (2012), \href{https://arxiv.org/abs/1003.0535}{arXiv:1003.0535}.
\bibitem{gmw} C. G\'erard, M. Molodyk and M. Wrochna, \emph{Hadamard states for linearized gravity on asymptotically flat spacetimes}, \href{https://arxiv.org/abs/2610.31214}{arXiv:2610.31214}.
\bibitem{schmid} G. Schmid, \emph{On the Quantisation of Linear Gauge Theories on Lorentzian Manifolds: Maxwell's Theory via Complete Gauge Fixing}, \href{https://arxiv.org/abs/2604.27026}{arXiv:2604.27026}.
\bibitem{gow} C. G\'erard, O. Oulghazi and M. Wrochna, \emph{Hadamard states for the Klein--Gordon equation on Lorentzian manifolds of bounded geometry}, Comm. Math. Phys. 352 (2017), \href{https://arxiv.org/abs/1602.00930}{arXiv:1602.00930}.

\bibitem{p09}
\emph{Positive physical Hadamard states for critical Einstein--scalar multiplets on compact Cauchy surfaces}, archived manuscript 09 (2026), \href{https://github.com/g3173028019-beep/ai-automated-quantum-gravity-research/tree/main/papers/09-canonical-einstein-scalar-hadamard}{source and example}.
\bibitem{p21}
\emph{Coherent relative entropy of the Weyl algebra in a causal diamond}, archived manuscript 21 (2026), \href{https://github.com/g3173028019-beep/ai-automated-quantum-gravity-research/blob/main/papers/21-weyl-diamond-relative-entropy/manuscript.md}{manuscript}.
\bibitem{hack}
T.-P. Hack, \emph{Quantization of the linearised Einstein--Klein--Gordon system on arbitrary backgrounds and the special case of perturbations in inflation}, Class.\ Quantum Grav.\ \textbf{31}, 215004 (2014), \href{https://arxiv.org/abs/1403.3957}{arXiv:1403.3957}.
\bibitem{hackschenkel}
T.-P. Hack and A. Schenkel, \emph{Linear bosonic and fermionic quantum gauge theories on curved spacetimes}, \href{https://arxiv.org/abs/1205.3484}{arXiv:1205.3484}.
\bibitem{khavpeierls}
I. Khavkine, \emph{Covariant phase space, constraints, gauge and the Peierls formula}, Int.\ J.\ Mod.\ Phys.\ A \textbf{29}, 1430009 (2014), \href{https://arxiv.org/abs/1402.1282}{arXiv:1402.1282}.
\bibitem{khavlocal}
I. Khavkine, \emph{Local and gauge invariant observables in gravity}, Class.\ Quantum Grav.\ \textbf{32}, 185019 (2015), \href{https://arxiv.org/abs/1503.03754}{arXiv:1503.03754}.
\bibitem{petz}
D. Petz, \emph{Sufficient subalgebras and the relative entropy of states of a von Neumann algebra}, Commun.\ Math.\ Phys.\ \textbf{105}, 123--131 (1986), \href{https://math.bme.hu/~petz/pdf/29suff.pdf}{author copy}.
\bibitem{araki}
H. Araki, \emph{Von Neumann algebras of local observables for free scalar field}, J.\ Math.\ Phys.\ \textbf{5}, 1--13 (1964), \href{https://doi.org/10.1063/1.1704063}{doi:10.1063/1.1704063}.
\bibitem{longo}
R. Longo, \emph{Modular structure of the Weyl algebra}, Commun.\ Math.\ Phys.\ (2022), \href{https://arxiv.org/abs/2111.11266}{arXiv:2111.11266}.
\bibitem{longofinite}
R. Longo, V. Morinelli, F. Preta and K.-H. Rehren, \emph{Split property for free massless finite helicity fields}, Ann.\ Henri Poincar\'e \textbf{20}, 2555--2584 (2019), \href{https://arxiv.org/abs/1806.04967}{arXiv:1806.04967}.
\end{thebibliography}
\end{document}

